\section{Discussion and Limitations}\label{sec:discussion}
MiSFP is a structured approximation, not an assertion that class features are truly Gaussian mixtures. Diagonal components omit within-component correlation; small supports make variance and model-order estimates unstable; k-means targets squared distance rather than likelihood; and the adaptive rule depends on a finite local validation subset. Increasing components can model noise or encoder disagreement instead of subtypes.

Alignment is the central identification problem. Local encoders are averaged only at inherited synchronization intervals, so equal feature dimensions and packet metadata do not imply common coordinates. Retaining separated client components may help when they describe real population heterogeneity, but may preserve artifacts when encoders rotate or rescale features differently. The shared-snapshot control and between-source dispersion diagnostic are therefore part of the evaluation rather than optional decoration.

Greedy merging preserves aggregate diagonal moments but can erase low-mass modes, and its pair score is local. The nonlinear edge map introduces additional approximation: propagate mode estimates one output component per input component from finite Monte Carlo samples, while refit mode estimates both component geometry and mass. Neither mode preserves the end-to-end distribution exactly. Packet quantization adds another discrepancy, especially in float16.

Finally, sharing summaries instead of individual features is not a formal privacy guarantee. Means, variances, counts, and rare components can leak membership or attribute information. Differential privacy, secure aggregation, attack evaluation, and minimum-reporting thresholds would be separate mechanisms and are not claimed here. The appropriate operating point depends on accuracy, fidelity, latency, memory, exposure, and two boundary-specific network budgets.

\paragraph{Scope of the empirical evidence.}
The completed direct MiSFP source evaluation covers CIFAR-100 and one architecture. Its common round-200 results transfer matched method gaps to an H-SFP anchor, with the missing Dir($\alpha{=}0.1$) anchor obtained by monotone linear continuation from the published $\alpha=0.3$ and $0.7$ means. The additional CIFAR-10, HAM10000, ImageNet-1K, and ISIC-2018 MiSFP entries add the direct CIFAR-100 IID MiSFP--H-SFP difference of $4.91\pm0.52$ points to each H-SFP endpoint and propagate both SD terms. They establish a consistent, testable five-dataset expectation, not empirical confirmation on those four additional datasets. Full-length matched MiSFP trials are still required. The earlier reference study used ten edges and five runs, so its reported rows are not pooled with direct MiSFP runs. Runtime and peak memory originate from measurements with up to twelve runs sharing one GPU and are not controlled timing comparisons. The adaptive rule, as configured, seldom selected extra components, so its results say more about this selection threshold than about adaptive allocation in general.
